\section{Planning Results}
\label{sec:results}

\paragraph{Model comparison.}
Table~\ref{tab:headline} addresses RQ1. MLP JEPA with GAR attains $.433$ strict success. Its mean is close to the token LM at $.422$, below the sentence LM at $.471$, and below the sentence-plus-latent model at $.521$. Causal JEPA attains $.366$. MLP JEPA has the smallest standard deviation across the five tested seeds, although the sample is too small to support a general stability claim.

Non-generative latent consequence prediction supports closed-loop planning under the common interface, but it does not outperform the sentence models. The leading sentence-plus-latent row is consistent with complementary benefits from generative and predictive objectives.

\begin{table}[t]
\caption{Closed-loop planning on stylized iGSM. Values are mean $\pm$ sample standard deviation over five seeds. Every method scores the same shuffled current feasible-action menu.}
\label{tab:headline}
\centering
\small
\begin{tabular}{lcc}
\toprule
Model & Strict success & Success with $+2$ actions \\
\midrule
Sentence LM + latent & $\mathbf{.521\pm.047}$ & $.809\pm.020$ \\
Sentence LM & $.471\pm.046$ & $.813\pm.023$ \\
MLP JEPA + GAR & $.433\pm.019$ & $.796\pm.033$ \\
Token LM & $.422\pm.124$ & $\mathbf{.836\pm.073}$ \\
Causal JEPA + GAR & $.366\pm.047$ & $.778\pm.043$ \\
\bottomrule
\end{tabular}
\end{table}

\paragraph{Objective ablation.}
Table~\ref{tab:central-ablation} addresses RQ2. Factual and counterfactual latent prediction without GAR attains $.183$ strict success. Predictive GAR attains $.433$. Factual prediction without GAR or counterfactual prediction attains $.076$. Ranking without pairwise calibration attains $.358$. GAR without the counterfactual latent loss attains $.345$.

Same-state action ordering accounts for the larger part of the improvement. Counterfactual successor prediction supplies an additional gain beyond the order labels. Merely exposing the predictor to factual transitions is insufficient.

\begin{table}[t]
\caption{Five-seed objective ablations. Rank is the GAR pairwise hinge. Reg. calibrates energy differences to geometric target differences. CF pred. predicts counterfactual successor states.}
\label{tab:central-ablation}
\centering
\small
\begin{tabular}{lccccl}
\toprule
Configuration & Rank & Reg. & CF pred. & Strict & $+2$ \\
\midrule
Full predictive GAR & yes & yes & yes & $.433\pm.019$ & $.796\pm.033$ \\
GAR ranking only & yes & no & yes & $.358\pm.045$ & $.750\pm.018$ \\
GAR without CF prediction & yes & yes & no & $.345\pm.049$ & $.734\pm.017$ \\
Latent prediction only & no & no & yes & $.183\pm.016$ & $.628\pm.031$ \\
Factual prediction only & no & no & no & $.076\pm.013$ & $.422\pm.022$ \\
\bottomrule
\end{tabular}
\end{table}

The reference horizon produces a sharply non-monotonic result. Horizons one, two, and four attain $.161$, $.433$, and $.181$ strict success. Horizon two is best among the tested settings, but the experiment does not support a universal optimum. Increasing counterfactual breadth from one to four changes strict success from $.403$ to $.439$. Pairwise-regression weights $.10$ and $.50$ attain $.443$ and $.402$, so the gain persists across nearby coefficients.

\paragraph{Information controls.}
Shuffling geometric quality labels while retaining counterfactual transitions gives $.015$ strict success at each of three screened learning rates. Direct GAR without successor auxiliaries reaches $.150$ to $.185$. Adding the JEPA successor objectives to the same direct policy path raises one-seed performance to $.440$, compared with $.435$ for transition-value GAR in that seed.

Counterfactual volume alone cannot explain the GAR result. Predictive auxiliary training clearly helps a direct scorer, but the current in-distribution comparison does not show that the policy itself must factor through a predicted successor. The stronger factorization claim requires an advantage in data efficiency, compositional length, semantic robustness, goal transfer, recovery after errors, or valid multi-step simulation.
